\def\level{Première}  % Classe à droite
\def\course{NSI}
\def\chaptername{Les dictionnaires} % Titre au centre
\def\docname{LP07}        % Identifiant à gauche

\pagestyle{cours_FP}
\makeheader{} % Affiche le bandeau

\section{Introduction}

\begin{definition}[Les dictionnaires]{}
	Un \textbf{dictionnaire} est une structure de données qui permet de stocker des paires \textbf{clé-valeur}. Les clés sont uniques dans un dictionnaire, tandis que les valeurs peuvent être dupliquées.

	En Python, les dictionnaires:
	\begin{itemize}
		\item sont des collections \textbf{non ordonnées} d'éléments;
		\item sont \textbf{mutables}, les éléments d'un dictionnaire peuvent être modifiés après sa création;
		\item sont \textbf{indexés par des clés}, qui peuvent être de n'importe quel type de données immuable;
		\item sont délimités par des \textbf{accolades} \verb|{ }| et les paires clé-valeur sont séparées par des \textbf{deux-points} \verb|:|.
	\end{itemize}
	
	\begin{CodePythontexAlt}[Largeur=1\linewidth,Centre,PremLigne=1,Lignes=false,]{}
mon_dictionnaire = {'a': 1, 'b': 2, 'c': 3}
	\end{CodePythontexAlt}

	Dans cet exemple:
	\begin{itemize}
		\item les clés sont 'a', 'b' et 'c';
		\item les valeurs sont 1, 2 et 3;
		\item la clé 'a' est associée à la valeur 1;
		\item le couple clé-valeur 'b': 2 signifie que la clé 'b' est associée à la valeur 2;
	\end{itemize}
	
\end{definition}

\begin{propriete}[Les dictionnaires]{}
\begin{CodePythontexAlt}[Largeur=1\linewidth,Centre,PremLigne=1]{}
mon_dictionnaire = {'a': 1, 'b': 2, 'c': 3}
print(mon_dictionnaire['a'])  # Affiche 1
print(mon_dictionnaire['b'])  # Affiche 2
print(mon_dictionnaire['c'])  # Affiche 3
print(len(mon_dictionnaire))  # Affiche 3
		\end{CodePythontexAlt}
\end{propriete}

\section{Itérer sur un dictionnaire}


\begin{minipage}{0.31\linewidth}
	\begin{easybox}{Itérer sur les clés}{}
		\begin{CodePythontexAlt}[Largeur=1\linewidth,Centre,Lignes=false,PremLigne=1]{}
	for cle in dico.keys():
		print(dico[cle])
		\end{CodePythontexAlt}
	\end{easybox}
\end{minipage}\hfill
\begin{minipage}{0.33\linewidth}
	\begin{easybox}{Itérer sur les valeurs}{}
		\begin{CodePythontexAlt}[Largeur=1\linewidth,Centre,Lignes=false,PremLigne=1]{}
	for val in dico.values():
		print(val)
		\end{CodePythontexAlt}
	\end{easybox}
\end{minipage}\hfill
\begin{minipage}{0.36\linewidth}
	\begin{easybox}{Itérer sur les items}{}
		\begin{CodePythontexAlt}[Largeur=1\linewidth,Centre,Lignes=false,PremLigne=1]{}
	for cle, val in dico.items():
		print(cle, val)
		\end{CodePythontexAlt}
	\end{easybox}
\end{minipage}

\section{Création de dictionnaires}

\begin{easybox}{3 façons de créer un dictionnaire}{}
	\begin{CodePythontexAlt}[Largeur=1\linewidth,Centre,PremLigne=1]{}
mon_dico = {}					   # Dictionnaire vide
mon_dico = {'a': 1, 'b': 2, 'c': 3} # Création avec des paires clé-valeur
mon_dico['a'] = 4				   # Création par ajout de cle : valeur
\end{CodePythontexAlt}
\end{easybox}

\begin{easybox}{Création d'un dictionnaire par itération}{}
	\begin{CodePythontexAlt}[Largeur=1\linewidth,Centre,PremLigne=1]{}
mon_dico = {}
for i in range(5):
	mon_dico[i] = i**2
print(mon_dico) # Affiche {0: 0, 1: 1, 2: 4, 3: 9, 4: 16}
\end{CodePythontexAlt}
\end{easybox}

\begin{easybox}{Création d'un dictionnaire par compréhension}{}
	\begin{CodePythontexAlt}[Largeur=1\linewidth,Centre,PremLigne=1]{}
mon_dico = {i: i**2 for i in range(5)}
print(mon_dico) # Affiche {0: 0, 1: 1, 2: 4, 3: 9, 4: 16}
\end{CodePythontexAlt}
\end{easybox}

\begin{exercice_cours}[Création d'un dictionnaire]{}
 Ecrire une fonction \texttt{longueurs\_elts(L)} qui prend en argument une liste de chaînes de caractères \texttt{L} et qui renvoie un dictionnaire dont les clés sont les éléments de \texttt{L} et les valeurs sont les longueurs des éléments de la liste \texttt{L}.
 
 
 \begin{CodePythontexAlt}[Largeur=1\linewidth,Centre,PremLigne=1, Lignes=false,]{}
def longueurs_elts(L):
	# A compléter
	

	






assert longueurs_elts(['a', 'ab', 'abc']) == {'a': 1, 'ab': 2, 'abc': 3}
\end{CodePythontexAlt}
\end{exercice_cours}


\section{Appartenance d'un élément à un dictionnaire}

	\begin{CodePythontexAlt}[Largeur=1\linewidth,Centre,PremLigne=1]{}
mon_dico = {'a': 1, 'b': 2, 'c': 3}
print('a' in mon_dico.keys())  # Affiche True
print('d' in mon_dico.keys())  # Affiche False
print(1 in mon_dico.values())  # Affiche True
print(4 in mon_dico.values())  # Affiche False
print(('a', 1) in mon_dico.items())  # Affiche True
print(('d', 4) in mon_dico.items())  # Affiche False
\end{CodePythontexAlt}

\section{Suppression dans un dictionnaire}

	\begin{CodePythontexAlt}[Largeur=1\linewidth,Centre,PremLigne=1]{}
mon_dico = {'a': 1, 'b': 2, 'c': 3}
del mon_dico['a']  # Supprime la paire clé-valeur 'a': 1
print(mon_dico)  # Affiche {'b': 2, 'c': 3}
mon_dico.clear()  # Supprime toutes les paires clé-valeur du dictionnaire
print(mon_dico)  # Affiche {}
\end{CodePythontexAlt}

\pagebreak

\begin{exercice_cours}{\quad}{}

			\begin{enumerate}
				\item  Ecrire une fonction \verb|occurence(sequence)| qui renvoie un dictionnaire dont les clés sont les caractères présents dans \verb|sequence| et leurs valeurs le nombre de fois où elles apparaissent.

				\begin{CodePythontexAlt}[Largeur=1\linewidth,Centre,Lignes=false,PremLigne=1]{}
					def occurence(sequence)->dict:
						  









					assert occurence([4, 2, 3, 4, 1, 4, 3, 4]) == {4: 4, 2: 1, 3: 2, 1: 1}
					assert occurence("abracadrabra") == {'a': 5, 'b': 2, 'r': 3, 'c': 1, 'd': 1}
				\end{CodePythontexAlt}

				\item Ecrire une fonction \verb|plus_frequent(sequence)| qui renvoie le caractère le plus souvent présent dans la \verb|sequence|.
				
				\begin{CodePythontexAlt}[Largeur=1\linewidth,Centre,Lignes=false,PremLigne=1]{}
					def plus_frequent(sequence):
						



					





					assert occurence([4, 2, 3, 4, 1, 4, 3, 4]) == 4
					assert occurence("abracadrabra") == 'a'
				\end{CodePythontexAlt}

			\end{enumerate}


\end{exercice_cours}

	

\begin{exercice_cours}{\quad}{}

	On considère un dictionnaire donnant la moyenne de quelques étudiants. Partitionner ce dictionnaire en deux sous dictionnaires : un dictionnaire
	des élèves ayant la moyenne et un dictionnaire des élèves n'ayant pas la
	moyenne.
	

	\begin{CodePythontexAlt}[Largeur=1\linewidth,Centre,Lignes=false,PremLigne=1]{}
		def coupe_dico(dico):
			









		dico={"etudiant1": 15,"etudiant2": 9,"etudiant3": 12,"étudiant4": 5}
		print(coupe_dico(dico))
	\end{CodePythontexAlt}
\end{exercice_cours}

